\subsection{Morfismos localmente de presentación finita}
\label{subsection:IV.1.4-fr}

\begin{env}[1.4.1]
\label{IV.1.4.1-fr}
Dado un anillo $A$, diremos que una $A$-álgebra $B$ (conmutativa)
es \emph{de presentación finita} si es isomorfa al cociente
$A[T_1,\ldots,T_n]/\mathfrak a$ de un álgebra de polinomios sobre $A$ por un
ideal \emph{de tipo finito} $\mathfrak a$ de $A[T_1,\ldots,T_n]$. Para toda
$A$-álgebra (conmutativa) $A'$, $B\otimes_AA'$ es entonces una $A'$-álgebra de
presentación finita. por ser isomorfa al cociente de $A'[T_1,\ldots,T_n]$ por
la imagen canónica en este anillo de $\mathfrak a\otimes_AA'$, que es
evidentemente un $A'[T_1,\ldots,T_n]$-módulo de tipo finito. Si $B$ es una
$A$-álgebra de presentación finita y $C$ una $B$-álgebra de presentación
finita, entonces $C$ es una $A$-álgebra de presentación finita. En efecto, sean
$B=A[S_1,\ldots,S_m]/\mathfrak a$ y
$C=B[T_1,\ldots,T_n]/\mathfrak b$, donde $\mathfrak a$ (resp. $\mathfrak b$) es
un ideal de tipo finito de $A[S_1,\ldots,S_m]$ (resp.
$B[T_1,\ldots,T_n]$); el anillo $B[T_1,\ldots,T_n]$ es isomorfo a
$A[S_1,\ldots,S_m,T_1,\ldots,T_n]/\mathfrak a'$, donde $\mathfrak a'$ es
la imagen canónica de
$\mathfrak a\otimes_AA[T_1,\ldots,T_n]$, y es por tanto un ideal de tipo finito de
$A[S_1,\ldots,S_m,T_1,\ldots,T_n]$; el ideal $\mathfrak b$ de
$B[T_1,\ldots,T_n]$ es de la forma $\mathfrak b'/\mathfrak a'$, donde
$\mathfrak b'$ es un ideal de $A[S_1,\ldots,S_m,T_1,\ldots,T_n]$; como
$\mathfrak a'$ y $\mathfrak b'/\mathfrak a'$ son módulos de tipo finito
sobre $A[S_1,\ldots,S_m,T_1,\ldots,T_n]$, lo mismo ocurre con $\mathfrak b'$,
y $C$, isomorfa a
$A[S_1,\ldots,S_m,T_1,\ldots,T_n]/\mathfrak b'$, es pues efectivamente una
$A$-álgebra de presentación finita. Observemos, por último, que si $A$ es
\emph{noetheriano}, lo mismo ocurre con $A[T_1,\ldots,T_n]$, por lo que toda
$A$-álgebra de tipo finito es entonces de presentación finita.
\end{env}

\begin{definition}[1.4.2]
\label{IV.1.4.2-fr}
Sean $f:X\to Y$ un morfismo de preesquemas, $x$ un punto de $X$, $y=f(x)$.
Se dice que $f$ es de presentación finita en $x$ si existe una vecindad
abierta afín $V$ de $y$ y una vecindad abierta afín $U$ de $x$ tales que
$f(U)\subset V$ y que el anillo $A(U)$ sea una $A(V)$-álgebra de presentación
finita. Se dice que $f$ es localmente de presentación finita si es de
presentación finita en todo punto de $X$.
\end{definition}

Si $Y$ es localmente noetheriano, es lo mismo decir que $f$ es
localmente de tipo finito o localmente de presentación finita.

\begin{proposition}[1.4.3]
\label{IV.1.4.3-fr}
\begin{enumerate}
  \item[\textup{(i)}] Todo isomorfismo local es localmente de presentación finita.
  \item[\textup{(ii)}] Si dos morfismos $f:X\to Y$, $g:Y\to Z$ son localmente de
  presentación finita, lo mismo ocurre con $g\circ f$.

\oldpage[IV-1]{231}

  \item[\textup{(iii)}] Si $f:X\to Y$ es un $S$-morfismo localmente de presentación
  finita, $f_{(S')}:X_{(S')}\to Y_{(S')}$ es localmente de presentación finita
  para toda extensión $S'\to S$ del preesquema de base.
  \item[\textup{(iv)}] Si $f:X\to X'$ y $g:Y\to Y'$ son dos $S$-morfismos
  localmente de presentación finita, $f\times_Sg$ es localmente de
  presentación finita.
  \item[\textup{(v)}] Si el compuesto $g\circ f$ de dos morfismos es localmente de
  presentación finita y si $g$ es localmente de tipo finito, entonces $f$ es
  localmente de presentación finita.
\end{enumerate}
\end{proposition}

\begin{proof}
La afirmación (i) es trivial. Para demostrar (iii), se puede limitar al caso
en que $S=Y$ (\textbf{I}, 3.3.11); sean $x'$ un punto de $X_{(S')}$, $s'$ y $x$ sus
proyecciones sobre $S'$ y $X$ respectivamente, $s$ la proyección común de $s'$
y $x$ sobre $S$. Por hipótesis, hay una vecindad abierta afín $V$ (resp.
$U$) de $s$ (resp. $x$) tales que la imagen de $U$ está contenida en $V$ y que
$A(U)$ es una $A(V)$-álgebra de presentación finita; sea $V'$ una vecindad
abierta afín de $s'$ cuya imagen en $S$ esté contenida en $V$; entonces
$U\times_VV'$ es una vecindad abierta afín de $x'$ cuya imagen en $S'$
está contenida en $V'$ y cuyo anillo $A(U)\otimes_{A(V)}A(V')$ es una
$A(V')$-álgebra de presentación finita (1.4.1).

Para demostrar (ii), consideremos un punto $x\in X$, y sean $y=f(x)$,
$z=g(f(x))$. Hay por hipótesis una vecindad abierta afín $W=\operatorname{Spec}(C)$
(resp. $V=\operatorname{Spec}(B)$) de $z$ (resp. $y$) tales que $g(V)\subset W$ y que $B$
es una $C$-álgebra de presentación finita. Por otra parte, hay una vecindad
abierta afín $V'=\operatorname{Spec}(B')$ (resp. $U=\operatorname{Spec}(A)$) de $y$ (resp. $x$) tales que
$f(U)\subset V'$ y que $A$ es una $B'$-álgebra de presentación finita. Sea
$s\in B$ tal que el abierto $D(s)=\operatorname{Spec}(B_s)$ en $V$ sea una vecindad de $y$
contenida en $V'$, y sea $U'=f^{-1}(D(s))\subset U$; se tiene
$U'=\operatorname{Spec}(A\otimes_{B'}B_s)$, y $A\otimes_{B'}B_s$ es una $B_s$-álgebra de
presentación finita (1.4.1); por otra parte $B_s=B[T]/(1\setmin sT)$ es una
$B$-álgebra de presentación finita por definición; por consiguiente (1.4.1),
$A\otimes_{B'}B_s$ es una $C$-álgebra de presentación finita, lo que prueba
(ii).

Se sabe que (iv) resulta de (i), (ii) y (iii) (\textbf{I}, 3.5.1). Para probar (v),
consideremos el diagrama conmutativo
\[
\xymatrix{
X \ar[r]^{\Gamma_f} \ar[d] & X\times_ZY \ar[d]^{f\times 1} \\
Y \ar[r]_{\Delta_g} & Y\times_ZY
}
\]
que identifica $X$ con el producto de los $(Y\times_ZY)$-preesquemas $Y$ y
$X\times_ZY$ (\textbf{I}, 5.3.7, siendo $\Delta_g$ el morfismo diagonal). Si se sabe
que $\Delta_g$ es localmente de presentación finita, resultará de
(iii) que lo mismo ocurre con $\Gamma_f$. Por otra parte, se tiene la factorización
de $f$ en
$X\xrightarrow{\Gamma_f}X\times_ZY\xrightarrow{p_2}Y$ (\textbf{I}, 5.3.13), donde la
proyección $p_2$ es igual a $(g\circ f)\times_Z1_Y$; como se supone que $g\circ f$ es
localmente de presentación finita, lo mismo ocurre con $p_2$ por
(iv), y resultará finalmente de (ii) que $f$ es también localmente de
presentación finita. Se ve, por tanto, que se reduce a demostrar el
\end{proof}

\begin{corollary}[1.4.3.1]
\label{IV.1.4.3.1-fr}
Sea $g:Y\to Z$ un morfismo localmente de tipo finito; entonces el morfismo
diagonal $\Delta_g:Y\to Y\times_ZY$ es localmente de presentación finita.
\end{corollary}

\oldpage[IV-1]{232}

\begin{proof}
Se puede limitar al caso en que $Z=\operatorname{Spec}(A)$, $Y=\operatorname{Spec}(B)$, siendo $B$ una
$A$-álgebra de tipo finito; el morfismo $\Delta_g$ corresponde entonces al
\emph{homomorfismo de aumentación} $\delta:B\otimes_AB\to B$ tal que
$\delta(x\otimes y)=xy$. Teniendo en cuenta la definición (1.4.2), basta
mostrar que el núcleo $\mathfrak J$ de $\delta$ es un ideal \emph{de tipo
finito}, lo que resulta de (\textbf{0}, 20.4.4).
\end{proof}

\begin{proposition}[1.4.4]
\label{IV.1.4.4-fr}
Sean $A$ un anillo, $B$ una $A$-álgebra, $B'$ un álgebra de polinomios
$A[T_1,\ldots,T_n]$, $u:B'\to B$ un homomorfismo sobreyectivo de
$A$-álgebras. Para que $B$ sea una $A$-álgebra de presentación finita, es necesario
y suficiente que el núcleo $\mathfrak J$ de $u$ sea un ideal de tipo finito de
$B'$.
\end{proposition}

\begin{proof}
La condición es suficiente por definición (1.4.1). Para ver que es
necesaria, observemos que el morfismo $g:\operatorname{Spec}(B')\to\operatorname{Spec}(A)$ es
localmente de tipo finito; si $B$ es una $A$-álgebra de presentación finita,
el morfismo $f:\operatorname{Spec}(B)\to\operatorname{Spec}(B')$ correspondiente a $u$ es tal que
$g\circ f:\operatorname{Spec}(B)\to\operatorname{Spec}(A)$ es localmente de presentación finita, por lo que
resulta de (1.4.3, (v)) que $f$ es también localmente de presentación finita.
Ahora bien, para demostrar que el ideal $\mathfrak J$ es de tipo finito, basta
probar que $\widetilde{\mathfrak J}$ es un $\widetilde{B}'$-Módulo de tipo
finito (\textbf{II}, 6.1.4.1). Volviendo a la definición (1.4.2), se reduce finalmente
a probar el
\end{proof}

\begin{lemma}[1.4.4.1]
\label{IV.1.4.4.1-fr}
Sean $\mathfrak a$ un ideal de un anillo $C$, $s$ un elemento de $C$, tal que
$C_s/\mathfrak a_s$ sea una $C$-álgebra de presentación finita; entonces
$\mathfrak a_s$ es un ideal de tipo finito en el anillo $C_s$.
\end{lemma}

\begin{proof}
Por hipótesis, existe una $C$-álgebra de polinomios
$C'=C[T_1,\ldots,T_n]$ y un $C$-homomorfismo sobreyectivo
$v:C'\to C_s/\mathfrak a_s$ cuyo núcleo $\mathfrak b$ es de tipo finito.
Sea $p:C_s\to C_s/\mathfrak a_s$ el homomorfismo canónico; para todo $i$,
hay un $t_i\in C$ tal que $p(t_i/s^k)=v(T_i)$ (se puede suponer que
el exponente $k$ es el mismo para todos los índices $i$). Consideremos entonces la
$C_s$-álgebra $D=C_s[T_1,\ldots,T_n]$, y el homomorfismo de $C_s$-álgebras
$w:D\to C_s$ tal que $w(T_i)=t_i/s^k$; $w$ es evidentemente sobreyectivo, y lo
mismo ocurre, por consiguiente, con $v'=p\circ w:D\to C_s/\mathfrak a_s$. Por otra
parte, todo polinomio de $D$ puede escribirse $P/s^m$, donde
$P\in C'=C[T_1,\ldots,T_n]$, y se tiene
$v'(P/s^m)=(1/s^m)v(P)$; como $1/s$ es invertible en $C_s$, la relación
$v'(P/s^m)=0$ en $C_s/\mathfrak a_s$ equivale a $v(P)=0$, y por consiguiente el
núcleo $\mathfrak b'$ de $v'$ está generado por la imagen canónica de
$\mathfrak b$ en $D$, y \emph{a fortiori} es \emph{de tipo finito} en el anillo
$D$. Pero
$\mathfrak b'=v'^{-1}(0)=w^{-1}(p^{-1}(0))=w^{-1}(\mathfrak a_s)$, y como
$w$ es sobreyectivo, $\mathfrak a_s=w(\mathfrak b')$, lo que prueba que
$\mathfrak a_s$ es un ideal de tipo finito de $C_s$.
\end{proof}

\begin{corollary}[1.4.5]
\label{IV.1.4.5-fr}
Sean $X$, $Y$ dos preesquemas, $j:X\to Y$ una inmersión, $U$ un abierto de
$Y$ tal que $j(X)$ sea cerrado en $U$, $\mathcal J$ el Ideal cuasicoherente de
$\mathcal O_U$ que define el subpreesquema cerrado de $Y$ asociado a $j$. Para
que $j$ sea localmente de presentación finita, es necesario y suficiente que
$\mathcal J$ sea un $\mathcal O_U$-Módulo de tipo finito.
\end{corollary}

\par\medskip\noindent
\begin{proposition}[1.4.6]
\label{IV.1.4.6-fr}
Sea $\varphi:A\to B$ un homomorfismo de anillos. Para que el morfismo
correspondiente $f:\operatorname{Spec}(B)\to\operatorname{Spec}(A)$ sea localmente de presentación finita,
es necesario y suficiente que $B$ sea una $A$-álgebra de presentación finita.
\end{proposition}

\begin{proof}
Siendo la condición trivialmente suficiente, queda por probar que es
necesaria. Si $f$ es localmente de presentación finita, resulta ya de
(1.3.6) que $B$ es una $A$-álgebra de tipo finito, por lo que existe un
homomorfismo sobreyectivo $u:B'=A[T_1,\ldots,T_n]\to B$ de $A$-álgebras;
resulta de (1.4.3, (v)) que la inmersión
$j:\operatorname{Spec}(B)\to\operatorname{Spec}(B')$ es localmente
\oldpage[IV-1]{233}
de presentación finita; por tanto (1.4.5), si $\mathfrak J$ es el núcleo de $u$,
el $\widetilde{B}'$-Módulo $\widetilde{\mathfrak J}$ es de tipo finito, y por
consiguiente (\textbf{II}, 6.1.4.1) $\mathfrak J$ es un ideal de tipo finito.
\end{proof}

\begin{proposition}[1.4.7]
\label{IV.1.4.7-fr}
Sea $\rho:A\to B$ un homomorfismo de anillos que hace de $B$ una $A$-álgebra
finita. Para que $B$ sea una $A$-álgebra de presentación finita, es necesario y
suficiente que $B$ sea un $A$-módulo de presentación finita.
\end{proposition}

Utilizaremos el

\begin{lemma}[1.4.7.1]
\label{IV.1.4.7.1-fr}
Si $B$ es una $A$-álgebra finita, existe un $A$-homomorfismo sobreyectivo
de álgebras $u:B'\to B$, donde $B'$ es una $A$-álgebra finita y de presentación
finita que es un $A$-módulo libre.
\end{lemma}

\begin{proof}
En efecto, hay un sistema finito $(a_i)_{1\leq i\leq m}$ de generadores de la
$A$-álgebra $B$ tal que cada $a_i$ verifica una relación $F_i(a_i)=0$, donde
$F_i\in A[X]$ es un polinomio mónico de grado $>0$; la $A$-álgebra cociente
$B'_i=A[X]/(F_i)$ es, por tanto, libre de rango finito sobre $A$; sea $c_i$ la clase de
$X$ en $B'_i$. Existe un $A$-homomorfismo de álgebras $u_i:B'_i\to B$
tal que $u_i(c_i)=a_i$; basta entonces tomar como $B'$ el producto
tensorial $B'_1\otimes_AB'_2\otimes\ldots\otimes_AB'_m$ y considerar
el homomorfismo $u_1\otimes u_2\otimes\ldots\otimes u_m:B'\to B$, que es
sobreyectivo por construcción. Además, si $B''=A[T_1,\ldots,T_m]$ y si $b'_i$
denota la imagen canónica de $c_i$ en $B'$, el $A$-homomorfismo de álgebras
$v:B''\to B'$ tal que $v(T_i)=b'_i$ ($1\leq i\leq m$) es sobreyectivo, y su
núcleo $\mathfrak b'$ es el ideal de $B''$ generado por los $F_i(T_i)$, por tanto de
tipo finito; esto prueba que $B'$ es una $A$-álgebra de presentación finita.
\end{proof}

\begin{proof}
Una vez demostrado este lema, conservemos las mismas notaciones, y sea
$w=v\circ u:B''\to B$. Si $\mathfrak b$ (resp. $\mathfrak a$) es el núcleo de
$w$ (resp. $u$), se tiene $\mathfrak a=v(\mathfrak b)$ puesto que $v$ es sobreyectivo.
Ahora bien (1.4.4), si $B$ es una $A$-álgebra de presentación finita, $\mathfrak b$
es un ideal de tipo finito de $B''$, por lo que $\mathfrak a$ es un ideal de tipo
finito de $B'$, y por tanto un $A$-módulo de tipo finito puesto que $B$ es una $A$-álgebra
finita; como $B'$ es un $A$-módulo libre, $B$ es un $A$-módulo de
presentación finita. Recíprocamente, si $B'$ es un $A$-módulo de presentación
finita, $\mathfrak a$ es un $A$-módulo de tipo finito (Bourbaki,
\emph{Alg. comm.}, cap. I, \S~2, n\textsuperscript{o}~8, lema 9) y
\emph{a fortiori} un ideal de tipo finito de $B'$; por consiguiente, $B$ es por
definición una $B'$-álgebra de presentación finita, y como $B'$ es una
$A$-álgebra de presentación finita, $B$ es una $A$-álgebra de presentación
finita.
\end{proof}

